\documentclass[11pt]{article}
\usepackage[margin=1in]{geometry}
\usepackage{amsmath,amssymb,amsthm,mathtools,microtype}
\usepackage[hidelinks]{hyperref}
\newtheorem{theorem}{Theorem}
\newtheorem{lemma}[theorem]{Lemma}
\newcommand{\E}{\mathbb E}
\newcommand{\Var}{\operatorname{Var}}
\newcommand{\sgn}{\operatorname{sgn}}
\newcommand{\diag}{\operatorname{diag}}
\newcommand{\wh}{\widehat}
\title{A short parallel-query proof for fourfold forrelation}
\author{Amir H. Safavi-Naeini}
\date{}
\begin{document}
\maketitle
\begin{abstract}
A single parallel quantum interrogation of four binary phase masks needs
hard photon dose at least $(2/15)N^{1/8}$ to decide fourfold forrelation
with error at most $1/3$, for powers of two $N\ge2^{30}$.
The proof has three parts: a matrix bound on Fourier coefficients,
a correlated sign distribution, and concentration onto the exact promise.
Arbitrary idlers, repeated modes, and coherence between photon-number
sectors are allowed. Outcome-dependent fresh batches require an
additional Fourier bound.
\end{abstract}

\section{Statement and proof outline}
Let $N=2^d$, with coordinates in $\mathbb F_2^d$, and set
\[
 H_{ij}=N^{-1/2}(-1)^{i\cdot j},\qquad
 u=N^{-1/2}(1,\ldots,1)^{\mathsf T}.
\]
For four sign vectors $x^{(b)}\in\{\pm1\}^N$ and
$D_b=\diag(x^{(b)})$, define
\[
 F(x)=\operatorname{Forr}_4(x)=u^{\mathsf T}D_1HD_2HD_3HD_4u .
\]
The task is to distinguish $F\ge1/4$ from $F\le-1/4$.
A parallel probe is prepared before any sample output is available,
passes through the direct-sum mask $D_1\oplus\cdots\oplus D_4$ once,
and is measured arbitrarily. Hard dose $D$ bounds the total number of
signal photons on the entire support of the prepared state.
References and idlers bypass the masks and are free.

\begin{theorem}\label{thm:main}
For every power of two $N\ge2^{30}$, a parallel probe solving the task
with worst-case error at most $1/3$ has
\[
                     D\ge\frac{2}{15}N^{1/8}.
\]
This includes mixtures of probes and predetermined collections of batches
whose total hard dose is $D$: they can be prepared jointly and measured
by a single final POVM.
\end{theorem}

We construct input laws $\mu_+,\mu_-$ with mean forrelation approximately
$\pm(2\rho/\pi)^3$, where $\rho=200/201$. Their low-order sign correlations
are small. Lemma~\ref{lem:fourier} limits how strongly a parallel
measurement can collect those correlations. Concentration then puts the
mass of each law predominantly on its required promise.

The proof below uses standard Gaussian Hermite expansion, the Gaussian
arcsine identity, and the Efron--Stein and Gaussian Poincare inequalities.
The arithmetic check \path{verify_constants.py} verifies the final
constants. See \path{../AUDIT.md} for the proof audit.

\section{The parallel Fourier bound}
Write $m=4N$ for the total number of sign coordinates, and
$\chi_S(x)=\prod_{a\in S}x_a$.
A bosonic probe of hard dose $D$ embeds in $D$ labelled slots with alphabet
$\{0,1,\ldots,m\}$, where $0$ denotes vacuum. Send each normalized
occupation basis vector with $n\le D$ photons to the normalized symmetric
sum of distinct slot words with those labels and $D-n$ vacuum symbols.
These sums are orthonormal, giving an isometry from the truncated Fock
space. The query is exactly $O_x^{\otimes D}$, with
$O_x|0\rangle=|0\rangle$ and $O_x|a\rangle=x_a|a\rangle$.
Enlarging to arbitrary slot states can only strengthen the competitor.

For a word $a=(a_1,\ldots,a_D)$ let $\partial a$ be the set of nonzero
labels occurring oddly. Purify the input and write
$\psi=\sum_a|a\rangle\otimes\psi_a$, where the $\psi_a$ are idler vectors.
For an acceptance effect $0\le E\le I$, put
$E_{ab}=(\langle a|\otimes I)E(|b\rangle\otimes I)$. Then
\[
 f(x)=\sum_{a,b}\chi_{\partial a\triangle\partial b}(x)
                  \langle\psi_a,E_{ab}\psi_b\rangle,
 \qquad \sum_a\|\psi_a\|^2=1.
 \tag{1}
\]
Thus $\wh f(S)=0$ for $|S|>2D$. Arbitrary idlers are included in the
operator blocks $E_{ab}$. The Fourier coefficient is
$\wh f(S)=2^{-m}\sum_x f(x)\chi_S(x)$.

We use the elementary Schur factorization bound. If
$K_{ab}=\langle p_a,q_b\rangle$ with $\|p_a\|\le A$ and $\|q_b\|\le B$,
then
\[
 \|[K_{ab}E_{ab}]\|\le AB\|E\|.
 \tag{2}
\]
Map $(v_b)_b$ to $(q_b\otimes v_b)_b$, apply $I\otimes E$, and contract
with the $p_a$ to obtain (2). Let $\gamma_2(K)$ be the infimum of $AB$.

\begin{lemma}\label{lem:fourier}
For arbitrary complex numbers $c_S$ at level $1\le r\le2D$,
\[
 \begin{aligned}
 \left|\sum_{|S|=r}c_S\wh f(S)\right|
 &\le \frac12\max_{|S|=r}|c_S|\sum_{k=0}^r {D\choose k}{D\choose r-k}
                  m^{\min(k,r-k)/2}\\
 &\le \frac12{2D\choose r}m^{r/4}\max_{|S|=r}|c_S|.
 \end{aligned}
 \tag{3}
\]
\end{lemma}
\begin{proof}
In a pair of words $(a,b)$ with $|\partial a\triangle\partial b|=r$,
each differing label occurs oddly on exactly one side.
Mark its last occurrence on that side. This gives a unique record
of $r$ marked positions among the $D$ ket and $D$ bra slots.

Fix a record with $k$ ket marks. The marked labels are tuples
$\alpha(a)$ of length $k$ and $\beta(b)$ of length $r-k$.
Let $A,B$ be their underlying sets and put
\[
 S_a=\partial a\triangle A,\qquad S_b=\partial b\triangle B .
\]
A word is admissible on its own side if its marked labels are nonzero
and distinct, each mark is the last occurrence of its label, and its
residual parity is disjoint from its marked set. Write these two
indicators as $\ell(a)$ and $r(b)$. Define
\[
 C_{\alpha\beta}=
 \begin{cases}
 c_{A\cup B},&\text{all $r$ marked labels are distinct},\\
 0,&\text{otherwise}.
 \end{cases}
\]
The exact record contribution to (1) is the Schur symbol
\[
 K^{R}_{ab}=\ell(a)r(b)\,
             C_{\alpha(a),\beta(b)}\,\mathbf1_{\{S_a=S_b\}} .
 \tag{4}
\]
Residual equality cancels every unmarked parity. The two disjointness
conditions put each marked label on its specified odd side, and the
last-occurrence conditions make the record unique.
Repeated labels and vacuum slots are retained.

A factorization of $C$ lifts to (4) by tensoring its row and column
factors with $|S_a\rangle$ and $|S_b\rangle$, then multiplying by
the indicators. Thus $\gamma_2(K^R)\le\gamma_2(C)$.
The matrix $C$ has at most $m^k$ rows and $m^{r-k}$ columns.
Standard basis vectors on its shorter side give
\[
 \gamma_2(C)\le m^{\min(k,r-k)/2}\max|c_S|.
\]
Since $r>0$, $K^R$ has zero diagonal: at $a=b$, residual equality forces
$A=B$, while distinctness forces $A\cap B=\varnothing$, contradicting
$|A|+|B|=r$. Hence
\[
 [K^R_{ab}E_{ab}]=[K^R_{ab}(E-I/2)_{ab}],\qquad
 \|E-I/2\|\le1/2.
\]
Apply (2), and sum over the $ {D\choose k}{D\choose r-k}$ records.
Vandermonde's identity finishes the proof.
\end{proof}

\section{Correlated signs with small moments}
For $s=1,2,3$, independently draw standard Gaussian vectors $Z_s,W_s$
with cross-covariance $\E Z_sW_s^{\mathsf T}=\rho H$.
This is a valid covariance because $H$ is orthogonal and $\rho<1$.
Set $a_s=\sgn Z_s$, $b_s=\sgn W_s$, applying $\sgn$ componentwise, and form
\[
 y=(a_1,\ b_1\odot a_2,\ b_2\odot a_3,\ b_3).
\]
Here $\odot$ denotes componentwise multiplication.
Choose three independent uniform signs $\zeta_1,\zeta_2,\zeta_3$,
put $\zeta_4=\sigma\zeta_1\zeta_2\zeta_3$, and set
$x^{(b)}=\zeta_b y^{(b)}$. This defines $\mu_\sigma$.

For one link let
$L(A,B)=\E\prod_{i\in A}a_i\prod_{j\in B}b_j$.
If $a=|A|$, $b=|B|$ and $a+b\le r\le\sqrt N$, then
\[
 |L(A,B)|\le\left(\frac{r}{2\sqrt N}\right)^{\max(a,b)}.
 \tag{5}
\]
If one set is empty the moment is zero unless both are empty.
Otherwise the cross-covariance matrix $R=\rho H_{A,B}$ has norm at most
$\rho\sqrt{ab/N}\le r/(2\sqrt N)$.
The sign products have Hermite degrees at least $a$ and $b$.
Conditional expectation pairs only equal degrees $p\ge\max(a,b)$,
where it acts as $R^{\otimes p}$ on symmetric tensors and has norm
at most $\|R\|^p$. Cauchy--Schwarz and the unit $L^2$ norms give (5).

For $S=(S_1,\ldots,S_4)$ of total size $r$, gauge averaging kills the
difference of its two moments unless every $q_b=|S_b|$ is odd.
In that case
\[
 (\E_{\mu_+}-\E_{\mu_-})\chi_S
       =2L(S_1,S_2)L(S_2,S_3)L(S_3,S_4).
\]
Taking weights $(3/4,1/4)$, $(1/2,1/2)$ and $(1/4,3/4)$ in the
three maxima proves
\[
 \max(q_1,q_2)+\max(q_2,q_3)+\max(q_3,q_4)\ge3r/4.
\]
Consequently
\[
 |(\E_{\mu_+}-\E_{\mu_-})\chi_S|
 \le 2^{\,1-3r/4}r^{3r/4}N^{-3r/8}.
 \tag{6}
\]
Only even levels $r\ge4$ can contribute.

\section{Concentration onto the promises}
The arcsine identity gives $\E a_i b_j=\beta_NH_{ij}$, with
\[
 \beta_N=\frac{2\sqrt N}{\pi}\arcsin(\rho/\sqrt N).
\]
The three links are independent. The defining $N^{-5/2}$ path-sum
coefficient, three factors $N^{-1/2}$, and $N^4$ terms give
\[
 \E_{\mu_\sigma}F=\sigma\beta_N^3,\qquad
 \beta_N^3-\frac14>
 \left(\frac{1400}{2211}\right)^3-\frac14>\frac1{260}.
 \tag{7}
\]
Here $\arcsin z\ge z$ and $\pi<22/7$ suffice.

\begin{lemma}\label{lem:variance}
For either sign law, $\Var F<1540/N$.
\end{lemma}
\begin{proof}
Represent each link with independent standard Gaussians $G,\xi,\eta$:
\[
 a=\sgn(G+\xi/\sqrt{200}),\qquad
 b=\sgn(HG+\eta/\sqrt{200}).
\]
After standardization its correlation is $\rho H$.
Conditional on $G_1,G_2,G_3$, all coordinates of the four ungauged
blocks $y$ are independent. Primitive conditional means are
$\phi(G)$ and $\phi(HG)$, where
\[
 \phi(z)=\operatorname{erf}(10z),\qquad
 L=\|\phi'\|_\infty=20/\sqrt\pi,\qquad L^2<128.
\]
The function $\phi$ is applied componentwise to vectors.
The gauges multiply $F$ by the deterministic sign $\sigma$.

For an independent biased sign $Y$, Efron--Stein bounds its variance
contribution by
$(1-(\E Y)^2)(\partial_YF)^2\le(\partial_YF)^2$,
where $\partial_YF$ is the coefficient of $Y$.
Each endpoint block has total squared derivative $1/N$.
For block two the derivative is $\ell_jr_j$, with
\[
 \ell=HD_{a_1}u,\qquad
 r=HD_{b_2\odot a_3}HD_{b_3}u.
\]
These vectors are independent, $\E\ell_j^2=1/N$, and $\|r\|=1$.
The contribution is $1/N$ after unconditional averaging.
Block three is the same calculation with the ends reversed. Hence
\[
                  \E\Var(F\mid G_1,G_2,G_3)\le4/N.
 \tag{8}
\]

Put $M(g)=D_{\phi(g)}HD_{\phi(Hg)}$, a contraction. The conditional
mean is
\[
 \overline F=\sigma u^{\mathsf T}M(G_1)M(G_2)M(G_3)u.
\]
At one link write this as $p^{\mathsf T}M(G)q$, up to $\sigma$.
The partial-chain vectors $p,q$ are independent of each other and
$G$, with norms at most one and
$\E|p_i|^2,\E|q_i|^2\le1/N$.
To justify the last statement, set
$(P_sv)_i=v_{i+s}$ and $(D_sv)_i=(-1)^{s\cdot i}v_i$.
The Sylvester identities and oddness of $\phi$ give
\[
 M(P_sg)=P_sM(g),\qquad M(D_sg)=M(g)P_s.
\]
Replacing the leftmost Gaussian by $P_sG$ in a right chain, or by $D_sG$
in a left chain written with transposed factors, translates the vector
by $P_s$. Gaussian invariance and transitivity give equal coordinate
second moments, whose sum is at most one; the empty chain is $u$.

Differentiation, using $H=H^{\mathsf T}$, gives
\[
 \|\nabla_G\overline F\|^2\le2L^2\left(
 \|p\odot HD_{\phi(HG)}q\|^2+
 \|q\odot HD_{\phi(G)}p\|^2\right).
\]
Independence and the second-moment bounds make each expected squared
norm at most $1/N$. Gaussian Poincare over the three links gives
$\Var\overline F\le12L^2/N$. With (8),
$\Var F\le(4+12L^2)/N<1540/N$.
\end{proof}

\section{Closing the bound}
By (7) and Chebyshev, for $N\ge2^{30}$,
\[
 \varepsilon_\sigma:=\mu_\sigma\{\sigma F<1/4\}
 \le104104000/N<1/10.
 \tag{9}
\]
Pointwise correctness of an acceptance probability $f\in[0,1]$ gives
\[
 \E_{\mu_+}f\ge\tfrac23(1-\varepsilon_+),\qquad
 \E_{\mu_-}f\le\tfrac13(1-\varepsilon_-)+\varepsilon_-.
\]
Consequently
\[
 \E_{\mu_+}f-\E_{\mu_-}f
 \ge\frac13-\frac23(\varepsilon_++\varepsilon_-)>\frac15.
 \tag{10}
\]
Suppose $D<(2/15)N^{1/8}$. Then $2D<\sqrt N$.
Combine (3) and (6), recall $m=4N$, and use
$ {2D\choose r}\le(2eD/r)^r$:
\[
 \begin{split}
 |\E_{\mu_+}f-\E_{\mu_-}f|
 &\le\sum_{\substack{4\le r\le2D\\r\text{ even}}}
       \left(2^{3/4}eD\,r^{-1/4}N^{-1/8}\right)^r\\
 &\le\frac{q^4}{1-q^2}<\frac1{12},\qquad
 q=2^{1/4}eDN^{-1/8}<\frac12.
 \end{split}
\]
Here $r^{-1/4}\le4^{-1/4}$, $e<3$, and $2^{1/4}<5/4$.
Since $1/12<1/5$, this contradicts (10) and proves Theorem~\ref{thm:main}.

\section{What remains for adaptive batches}
For an outcome-dependent fresh protocol, acceptance has degree at most
$2D$: multiply batch probabilities along each classical history.
A sufficient remaining estimate is
\[
 \left|\sum_{|S|=r}c_S\wh f(S)\right|
 \le {2D\choose r}(4N)^{r/4}\max_{|S|=r}|c_S|
 \quad(1\le r\le2D).
 \tag{11}
\]
With (11), the same even-level calculation gives a gap at most
$2q^4/(1-q^2)<1/6<1/5$, proving the same dose bound.
This estimate remains unproved for adaptive batches; \path{../AUDIT.md}
describes the gap in the earlier argument. The theorem established here
concerns parallel interrogation.
\end{document}
